\chapter{Stand der Technik}\label{ch:sota}

Der Stand der Technik im Bereich Cache-bewusster Trie-Index-Strukturen
gliedert sich in sechs thematische Cluster (Such-Familie A--F) und einen
parallelen Allokator-Korpus (A01--A23). Insgesamt 33 Such-Algorithmus-Paper
und 21 Allokations-Paper wurden tieflektuert; die Ergebnisse sind als
\verb|algorithm_profiles/sota/| persistierte XML-Konfigurationen in der
\texttt{comdare-cache-engine} verfuegbar (30 Profile in V9.5).

\section{Trie-Familie (Cluster A)}

\textbf{P01 ART} \cite{leis2013art} fuehrt den Adaptive Radix Tree mit
adaptiver Knotengroesse Node4/16/48/256 ein.
\textbf{P02 HOT} \cite{binna2018hot} optimiert die Trie-Hoehe mittels
multi-byte Praefixen und bit-vector-Indizes.
\textbf{P03 Masstree} \cite{mao2012masstree} kombiniert Trie + B-Baum mit
Permutations-Index pro INode.
\textbf{P04 CoCo-Trie} \cite{boffa2024coco} verwendet succinct-encoded
compact pages.
\textbf{P05 START} \cite{fent2020start} ergaenzt ART um Self-Tuning der
Knoten-Konkretisierung.
P09 LOUDS (Jacobson 1989) ist die succinct-Trie-Foundation.
\textbf{P10 SuRF} \cite{zhang2018surf} nutzt LOUDS fuer Range-Filter.

\section{Hybrid + B+-Familie (Cluster B)}

\textbf{P06 B$^2$-Tree} \cite{schmeisser2022b2tree} hat zwei-stufiges
Fanout fuer L1+L2-Cache-Awareness.
\textbf{P07 Wormhole} \cite{wu2019wormhole} kombiniert Hash-Tabelle mit
linearer Praefix-Sicht.
P11 CSS-Tree (Rao/Ross 1999) ist die erste Cache-Sensitive Search Tree.
P12 CSB+-Tree (Rao/Ross 2000) entwickelt CSS zum B+-Tree weiter.
P13 Hankins/Patel 2003 untersucht den Effekt der Knotengroesse auf
Cache-Performance.

\section{Layout-Theorie (Cluster C)}

P14 Samuel/Pedersen/Bonnet 2005 macht CSB+-Tree prozessor-bewusst.
P15 Graefe/Larson 2001 untersucht B-Tree-Indizes vs. CPU-Caches.
P16 Bender/Demaine/Farach-Colton 2002 fuehrt cache-oblivious
Tree-Layout (van Emde Boas) ein.
P17 Bender et al. 2005 erweitert es zu Cache-Oblivious B-Trees.
P18 Saikkonen/Soisalon-Soininen 2008 + P19 (2016) entwickeln
Multi-Level + Layout-Invariant B-Trees.

\section{Prefetching (Cluster D + E)}

P20 B-Trees-Are-Back (Mueller/Benson/Leis 2025) demonstriert moderne
B-Tree-Engineering.
P21 Chen/Gibbons/Mowry 2001 fuehrt Prefetching B+-Trees ein.
P22 Chen et al. 2002 entwickelt fractal Prefetching.
P23 Khan 2010 macht Cache-Prefetching dynamisch adaptiv.
P24 Naderan-Tahan/Sarbazi-Azad 2016 baut adaptive Filter auf
Cache-Prefetcher.
P25 Mahling/Weisgut/Rabl 2025 (Fetch Me If You Can) kombiniert
Range-Scan mit Hot-Path-Caching.
P26 Zhang FGCS 2024 + P27 Zhang ASPLOS 2025 (Hierarchical Prefetching)
sind die juengsten Prefetching-Schemata.
P28 Kuehn DaMoN 2023 motiviert datenbasierte Cache-Optimierung.

\section{Synchronization + TUD-Habich (Cluster F)}

P08 ART of Practical Synchronization (Leis 2016) fuehrt
Optimistic-Lock-Coupling fuer ART ein.
P29 RCU (McKenney OLS 2001) und P30 Hazard Pointers (Michael 2004) sind
die kanonischen Reklamations-Schemata.
P31 Ungethuem/Habich 2017 + P32 Schmidt/Habich 2025 sind
TU-Dresden-spezifische Hardware-Optimierungs-Ueberblicke.
P33 VAMPIR-Poster (Berthold/Habich 2023) liefert die OTF2-Profiler-
Integration.

\section{Allokator-Cluster (A01--A23)}

Die 21 Allokator-Paper sind in der \texttt{Allokator\_Matrix.txt} der
\texttt{comdare-cache-engine} systematisch erfasst. Tier-1-Repraesentanten
sind Hoard (A01, ASPLOS 2000), Slab (A02, USENIX 1994),
Michael Lock-Free (A03, PLDI 2004), mimalloc (A04, MSR-TR-2019-18),
jemalloc (A05), TCMalloc (A06), snmalloc (A07, ISMM 2019),
NUMAlloc (A09, ISMM 2023). Tier-2/3 ergaenzt CAMA (A12, RTAS 2011),
StarMalloc (A13, OOPSLA 2024), Crystalline-Reclamation (A17, PLDI 2024)
und PIM-malloc (A16, arXiv 2025).
